\chapter{Kunihiko Kodaira}\label{ch1954kunihiko_kodaira:kunihiko-kodaira-fields-medal-1954-a-comprehensive-report}

\begin{center}\rule{0.5\linewidth}{0.5pt}\end{center}

\textbf{Prepared:} July 2026 \textbf{Subject:} Kunihiko Kodaira
(b. March 16, 1915, Tokyo, Japan; d. July 26, 1997)
\textbf{Position:} Professor, Institute for Advanced Study, Princeton
\textbf{Doctoral Advisor:} Shizuo Kakutani (Tokyo Imperial University);
additional studies with Kunizo Yosida (1938--1941)
\textbf{Fields Medal Citation:} \emph{``For his contributions to the
theory of harmonic integrals and its applications to the theory of
K{\"a}hler manifolds and algebraic varieties.''}

\begin{center}\rule{0.5\linewidth}{0.5pt}\end{center}

\begin{prereqs}
Complex analysis (holomorphic functions, Cauchy's theorem),
basic differential geometry (manifolds, tangent spaces, Riemannian metrics),
Riemannian geometry (Laplace--Beltrami operator, Hodge theory),
sheaf cohomology.
\end{prereqs}

\begin{objectives}
Explain the Hodge decomposition theorem for K{\"a}hler manifolds;
state the Kodaira vanishing theorem and its significance;
describe the embedding theorem for compact complex manifolds;
introduce Serre duality for compact complex manifolds;
explain the Kodaira--Spencer theory of deformations;
discuss applications to algebraic geometry.
\end{objectives}

\section{Problem History}

The \textbf{Hodge decomposition theorem} states that on a compact
Riemannian manifold $(M, g)$, every de Rham cohomology class has a
unique harmonic representative. This was proved by W. V. D. Hodge in
1941 for compact K{\"a}hler manifolds, but his proof was purely
topological and analytic, not explicitly using harmonic forms. The
explicit representation via harmonic integrals was missing, and the
\textbf{Hodge conjecture}---that every rational cohomology class on a
compact complex manifold is a linear combination of classes of algebraic
cycles---remained the deepest open problem in algebraic geometry.

In the 1930s and 1940s, Weil and de Rham proved that de Rham
cohomology is isomorphic to singular cohomology with real coefficients.
But there was no general understanding of how the harmonic forms depend
on the complex structure of a K{\"a}hler manifold. The relationship
between the complex geometry and the Riemannian geometry was poorly
understood, and the interplay between holomorphic forms and harmonic
forms was not systematically developed.

The \textbf{Kodaira vanishing theorem} \cite{kodaira1964}---that for a positive line
bundle $L$ on a compact complex manifold, $H^q(X, K_X \otimes L) = 0$
for $q > 0$---was expected but unproved. Serre had formulated the
conjecture in 1953, but no proof was known. This vanishing theorem was
crucial for the study of algebraic varieties because it provided a tool
for computing cohomology groups of line bundles and for understanding
when a complex manifold admits an embedding into projective space.

The challenge was to develop a general theory of harmonic forms on
complex manifolds that would simultaneously solve the representation
problem, yield vanishing theorems, and connect to the classification of
algebraic varieties. Kodaira faced the fundamental difficulty of
combining the analytical tools of partial differential equations with
the geometric insights of complex analysis and algebraic geometry.

Kodaira's work was deeply influenced by the Japanese mathematical
tradition of functional analysis, developed by von Neumann, Hermann
Weyl, and others at the Institute for Advanced Study. Kodaira studied at
Tokyo Imperial University under Shizuo Kakutani, and during 1938--1941
studied with Kunizo Yosida, who was a student of von Neumann. This
training in functional analysis, particularly the theory of Hilbert
spaces and self-adjoint operators, proved essential for Kodaira's
development of the theory of harmonic integrals. The spectral theory of
elliptic operators on compact manifolds provided the analytical
foundation for Kodaira's approach.

The central problem was to extend Hodge's theory---originally developed
for differential forms on smooth manifolds---to the setting of complex
manifolds, where the additional structure of the complex coordinates
provides new opportunities. On a K{\"a}hler manifold, the complex
structure $J$ and the Riemannian metric $g$ are compatible:
$g(JX, JY) = g(X, Y)$ for all tangent vectors $X, Y$. This
compatibility leads to the fundamental relation between the exterior
derivative $d$ and its complex analogues $\partial$ and $\bar{\partial}$.

A K{\"a}hler manifold is a complex manifold $M$ equipped with a
Hermitian metric $h$ whose associated fundamental form $\omega(X, Y) =
g(JX, Y)$ is closed, i.e., $d\omega = 0$. The closedness of $\omega$ is
the defining condition that makes the manifold K{\"a}hler. This
condition implies that the complex structure, the Riemannian metric, and
the symplectic form $\omega$ are all compatible, giving the manifold a
rich geometric structure that combines elements of complex geometry,
Riemannian geometry, and symplectic geometry.

The de Rham cohomology $H^p_{\text{dR}}(M)$ of a compact K{\"a}hler
manifold carries additional structure beyond what is available for
general smooth manifolds. The Hodge decomposition refines the cohomology
groups by decomposing them into $(p,q)$-types, where each $H^{p,q}(M)$
consists of cohomology classes represented by forms of type $(p,q)$.
This refinement is possible precisely because of the K{\"a}hler
condition.

\section{How Was It Solved}

Kodaira developed the theory of \textbf{harmonic integrals} on compact
complex manifolds in a series of papers (1949--1953), published in his
collected works. The key construction is the Hodge inner product on
complex-valued $p$-forms:
\[\langle \alpha, \beta \rangle = \int_M \alpha \wedge \overline{*}\beta\]
where $*$ is the Hodge star operator associated to a Hermitian metric
$h$ on $M$. The \textbf{harmonic projection} $H: L^2 \Omega^p(M) \to
\mathcal{H}^p(M)$ onto the space of harmonic forms
$\mathcal{H}^p(M) = \{\alpha : \Delta \alpha = 0\}$ is orthogonal with
respect to this inner product.

Kodaira proved the \textbf{Hodge decomposition theorem}:
$L^2 \Omega^p(M) = \operatorname{Im} d \oplus \operatorname{Im} d^*
\oplus \mathcal{H}^p(M)$, and every de Rham class $[\alpha] \in
H^p_{\text{dR}}(M)$ has a unique harmonic representative. This gives
the isomorphism $H^p_{\text{dR}}(M) \cong \mathcal{H}^p(M)$. The proof
uses the spectral theory of the Laplace operator $\Delta = d d^* + d^*
d$, which is a self-adjoint elliptic operator on the space of smooth
$p$-forms. The elliptic regularity theorem ensures that harmonic forms
are smooth, and the compactness of $M$ implies that $\mathcal{H}^p(M)$
is finite-dimensional.

For K{\"a}hler manifolds, the Laplacian satisfies $\Delta_d =
2\Delta_{\bar{\partial}}$, so $d$-harmonic forms decompose as
$\partial$-harmonic and $\bar{\partial}$-harmonic parts. This yields
the \textbf{Hodge diamond}:
\[h^{p,q} = \dim H^q(M, \Omega^p) = \dim \mathcal{H}^{p,q}(M)\]
and the \textbf{Hodge decomposition of cohomology}: $H^n(M,
\mathbb{C}) = \bigoplus_{p+q=n} H^{p,q}(M)$, where $H^{p,q}(M)
\cong \overline{H^{q,p}(M)}$. The decomposition is a consequence of the
K{\"a}hler identities, which relate the operators $d$, $\bar{\partial}$,
and their adjoints.

The Hodge numbers satisfy the symmetry $h^{p,q} = h^{q,p}$ (complex
conjugation) and the \textbf{Hodge--Riemann bilinear relations}: on
$H^{p,q}(M)$, the form
\[(-1)^{n(n-1)/2} i^{p-q} \int_M \alpha \wedge \bar{\alpha} \wedge
\omega^{n-p-q}\]
is positive definite, where $\omega$ is the K{\"a}hler form. These
relations severely restrict which Betti numbers can occur for K{\"a}hler
manifolds and are fundamental for the study of algebraic cycles. The
Hodge diamond for a K{\"a}hler surface is symmetric about both the
horizontal and vertical axes.

The \textbf{Kodaira vanishing theorem} \cite{kodaira1964} follows from the
Bochner--Kodaira--Nakano identity, which relates the Laplacians
$\Delta_{\bar{\partial}}$ and $\Delta_{\partial}$ on $(0,q)$-forms with
values in $L$ to curvature terms:
\[\Delta_{\bar{\partial}} = \nabla^* \nabla + \text{curvature terms}\]
When $L$ is positive, the curvature terms dominate, forcing $H^q(X,
K_X \otimes L) = 0$ for $q > 0$. The identity is derived by computing
the commutator $[\Delta_{\bar{\partial}}, \phi]$ for a smooth function
$\phi$ and using the curvature of the line bundle $L$.

The proof of the vanishing theorem proceeds by the maximum principle:
if $s$ is a harmonic $(0,q)$-form with values in $K_X \otimes L$ ($q >
0$), the Bochner--Kodaira--Nakano identity shows that
\[\int_M |\nabla s|^2 + \text{curvature terms} = 0\]
When $L$ is positive, the curvature terms are positive, forcing $s = 0$.
This ``vanishing by positivity'' is a general principle in K{\"a}hler
geometry. The method of proof is a powerful example of using curvature
to control cohomology.

The \textbf{Kodaira vanishing theorem} \cite{kodaira1964} implies the Kodaira embedding
theorem: for a positive line bundle $L$ on a compact complex manifold
$X$, the higher cohomology groups $H^q(X, L^m) = 0$ for $q > 0$ and
$m \gg 0$, and the complete linear system $|L^m|$ gives an embedding
$X \hookrightarrow \mathbb{P}^N$. This criterion settled the
fundamental question of when a complex manifold is algebraic. The
embedding is constructed by choosing a basis of holomorphic sections of
$L^m$ for $m$ sufficiently large, and the vanishing theorem ensures
that these sections separate points and tangent vectors.

Kodaira also established the \textbf{Serre duality} for compact
complex manifolds, which relates the cohomology groups $H^q(X,
\mathcal{F})$ and $H^{n-q}(X, \mathcal{F}^* \otimes K_X)$ for any
coherent analytic sheaf $\mathcal{F}$ on a compact complex manifold $X$
of dimension $n$. The duality is realized by the natural pairing:
\[\langle \alpha, \beta \rangle = \int_X \alpha \wedge \beta\]
where $\alpha \in H^q(X, \mathcal{F})$ and $\beta \in H^{n-q}(X,
\mathcal{F}^* \otimes K_X)$. This pairing is non-degenerate, giving an
isomorphism $H^q(X, \mathcal{F}) \cong H^{n-q}(X, \mathcal{F}^*
\otimes K_X)^*$. For $\mathcal{F} = \mathcal{O}_X$, this gives
$H^q(X, \mathcal{O}_X) \cong H^{n-q}(X, K_X)^*$, which relates the
geometric genus $p_g = h^{0,n}$ to the irregularity $q = h^{0,1}$.

The \textbf{Kodaira embedding theorem} states that a compact complex
manifold $X$ admits a holomorphic embedding into projective space
$\mathbb{P}^N$ if and only if $X$ carries a positive line bundle. The
proof uses the Kodaira vanishing theorem to show that for sufficiently
large $m$, $H^q(X, L^m) = 0$ for $q > 0$, and the map
$\phi_{|L^m|}: X \to \mathbb{P}^{N_m}$ is an embedding for $m$ large
enough. This result established the deep connection between complex
geometry and algebraic geometry. A compact complex manifold that admits
a positive line bundle is called a \textbf{projective manifold}, and the
embedding theorem provides a precise criterion for when this occurs.

Kodaira's work on \textbf{deformations of complex structures}---showing
that small deformations of a compact complex manifold are parametrized
by $H^1(X, \Theta_X)$ where $\Theta_X$ is the tangent sheaf---founded
deformation theory and led to the Kodaira--Spencer map and the theory of
moduli spaces. The Kodaira--Spencer theory provides the foundation for
understanding how complex manifolds vary in families, and has become a
central tool in modern algebraic geometry. Given a smooth family of
complex manifolds $\pi: \mathcal{X} \to B$ with central fiber $X_0 =
\pi^{-1}(b_0)$, the Kodaira--Spencer map assigns to each tangent vector
$v \in T_{b_0} B$ a cohomology class in $H^1(X_0, \Theta_{X_0})$ that
measures the infinitesimal change of the complex structure.

The \textbf{Bochner technique} is a fundamental tool in Kodaira's work.
Given a Hermitian vector bundle $E$ over a compact K{\"a}hler manifold,
the \textbf{Weitzenb{\"o}ck formula} expresses the Laplacian
$\Delta_E$ on $E$-valued forms in terms of the connection Laplacian
$\nabla^* \nabla$ and curvature terms:
\[\Delta_E = \nabla^* \nabla + \mathcal{R}\]
where $\mathcal{R}$ involves the curvature of $E$. For a line bundle
$L$ with positive curvature, the curvature terms in the Weitzenb{\"o}ck
formula are positive, which forces harmonic forms with values in $L$ to
vanish in certain degrees. This is the key mechanism behind the
Kodaira vanishing theorem.

The \textbf{elliptic regularity theorem} plays a crucial role in
Kodaira's theory. If $\Delta \alpha = 0$ for a distribution $\alpha$,
then $\alpha$ is smooth. This ensures that the harmonic representatives
of cohomology classes are smooth forms, which is essential for the
integration that defines the Hodge inner product. The theorem follows
from the fact that the Laplacian is an elliptic operator, and elliptic
operators have the property that their kernel consists of smooth
sections.

\section{Impact of This Discovery}

The Kodaira embedding theorem states: a compact complex manifold $X$ is
projective (holomorphically embeddable in $\mathbb{P}^N$) if and only
if it has a positive line bundle. The proof uses the Kodaira vanishing
theorem to show that for sufficiently large $m$, the higher cohomology
$H^q(X, L^m) = 0$ for $q > 0$, so the complete linear system $|L^m|$
gives an embedding. This criterion settled the fundamental question of
when a complex manifold is algebraic. The theorem has far-reaching
consequences for the classification of complex manifolds and for the
study of algebraic varieties.

The Hodge decomposition provides fundamental topological invariants:
$b_n(M) = \sum_{p+q=n} h^{p,q}$ (Betti numbers), and the Hodge numbers
satisfy $h^{p,q} = h^{q,p}$ (symmetry) and the \textbf{Hodge--Riemann
bilinear relations}. These constraints severely restrict which Betti
numbers can occur for K{\"a}hler manifolds. For example, the second
Betti number of a K{\"a}hler surface must satisfy $b_2 \geq 1$, and the
parity of $b_1$ must be even. These constraints are consequences of the
Hodge decomposition and the symmetries of the Hodge diamond.

The \textbf{Hodge conjecture}---that $H^{2p}(X, \mathbb{Q}) \cap
H^{p,p}(X)$ is generated by algebraic cycles---remains one of the seven
Clay Millennium Problems. Kodaira's work provides the framework: the
Hodge decomposition makes the conjecture precise. The conjecture
asserts that the middle-dimensional cohomology classes of type $(p,p)$
are exactly the classes that can be represented by algebraic cycles.
This connects the analytic theory of harmonic forms to the algebraic
geometry of subvarieties.

Kodaira's work on \textbf{deformations of complex structures}---showing
that small deformations of a compact complex manifold are parametrized
by $H^1(X, \Theta_X)$ where $\Theta_X$ is the tangent sheaf---founded
deformation theory and led to the Kodaira--Spencer map and the theory of
moduli spaces. The Kodaira--Spencer map $\rho: T_b B \to H^1(X_b,
\Theta_{X_b})$ describes how the complex structure varies in a family,
and its image measures the infinitesimal deformations that can be
realized.

The theory of \textbf{moduli spaces} of complex manifolds owes much to
Kodaira's work. The moduli space $\mathcal{M}_g$ of Riemann surfaces of
genus $g \geq 2$ is a complex orbifold of dimension $3g-3$, and its
tangent space at a point $[C]$ is isomorphic to $H^1(C, \Theta_C)$. The
 Kodaira--Spencer theory provides the local structure of $\mathcal{M}_g$
and is fundamental to the study of families of algebraic varieties. The
dimension formula $h^1(C, \Theta_C) = 3g-3$ follows from the
Riemann--Roch theorem and Serre duality.

The \textbf{Kodaira dimension} $\kappa(X)$, defined as the growth rate
of the plurigenera $P_m(X) = h^0(X, K_X^m)$, is a fundamental
birational invariant. The classification of surfaces by Enriques and
Kodaira, based on the Kodaira dimension, was the prototype for the
minimal model program in higher dimensions. Surfaces with $\kappa =
-\infty$ are rational or ruled, $\kappa = 0$ are K3, Enriques, Abelian,
or hyperelliptic, $\kappa = 1$ are properly elliptic, and $\kappa = 2$
are of general type. The Kodaira dimension provides a way to organize
complex manifolds according to their birational geometry.

Serre duality, as extended by Kodaira to compact complex manifolds, has
profound implications for the study of coherent sheaves and their
cohomology. For a line bundle $L$ on a compact complex manifold $X$,
Serre duality gives $H^q(X, L) \cong H^{n-q}(X, L^{-1} \otimes K_X)^*$,
which allows the computation of cohomology groups in terms of their
duals. This duality is essential for the Riemann--Roch theorem and for
the study of linear systems on algebraic varieties. The duality also
implies that the Euler characteristic
$\chi(X, L) = \sum_{q=0}^n (-1)^q h^q(X, L)$ can be computed from the
dual cohomology groups.

The \textbf{Riemann--Roch theorem} for line bundles on compact complex
manifolds, proved by Hirzebruch using the Hirzebruch--Riemann--Roch
theorem, expresses the Euler characteristic $\chi(X, L) = \sum_{q=0}^n
(-1)^q h^q(X, L)$ in terms of topological invariants:
\[\chi(X, L) = \int_X \operatorname{ch}(L) \cdot \operatorname{Td}(X)\]
where $\operatorname{ch}(L) = \exp(c_1(L))$ is the Chern character and
$\operatorname{Td}(X)$ is the Todd class of the tangent bundle. This
theorem generalizes the classical Riemann--Roch theorem for curves and
is a cornerstone of modern algebraic geometry. The
Hirzebruch--Riemann--Roch theorem uses the Atiyah--Singer index theorem
to compute the Euler characteristic in terms of characteristic classes.

The work of Kodaira on harmonic integrals also influenced the
development of \textbf{geometric analysis}, the study of partial
differential equations on manifolds. The elliptic regularity theory for
the Laplace operator, which Kodaira used to prove the smoothness of
harmonic representatives, became a model for the study of other
geometric PDEs. The Bochner technique, which Kodaira systematized,
remains a powerful tool in differential geometry. Many problems in
geometry can be reduced to the study of harmonic forms, and Kodaira's
methods provide the analytical tools for such investigations.

The \textbf{Hodge index theorem} for K{\"a}hler manifolds is a
fundamental consequence of the Hodge decomposition. It states that the
intersection form on $H^{1,1}(X) \cap H^2(X, \mathbb{R})$ has signature
$(1, h^{1,1}-1)$, where the positive direction is given by the
K{\"a}hler cone. This theorem is essential for the study of the
geometry of K{\"a}hler manifolds and has applications to the
classification of algebraic surfaces. The Hodge index theorem
generalizes the classical result for projective surfaces to higher
dimensions.

The theory of \textbf{harmonic maps} between K{\"a}hler manifolds,
developed by Eells and Sampson, extends Kodaira's techniques to the
study of maps that minimize the energy functional. Harmonic maps between
K{\"a}hler manifolds have special properties, including the fact that
holomorphic and antiholomorphic maps are harmonic. The study of harmonic
maps has led to deep results in differential geometry and topology,
including the proof of the sphere theorem and the study of negatively
curved manifolds.

\section{Modern Views and Improvements}

The Kodaira vanishing theorem has been extended in many directions.
Kawamata--Viehweg vanishing extends it to nef and big line bundles,
proving the Miyaoka--Yau inequality for surfaces. Nadel vanishing uses
multiplier ideals to obtain vanishing results for singular varieties.
The Nakano vanishing theorem generalizes to vector bundles using the
Bochner--Kodaira--Nakano identity. These generalizations have broad
applications in algebraic geometry, including the study of
higher-dimensional varieties and the minimal model program.

The Hodge decomposition has been generalized to non-compact K{\"a}hler
manifolds by Cattani--El Zeint--Kaplan and others, and to Hodge
modules by Saito. The \textbf{variation of Hodge structure}, studying
how $H^{p,q}$ varies with parameters, is a central topic in algebraic
geometry and has connections to mirror symmetry. The theory of \textbf{mixed
Hodge structures}, developed by Deligne, extends the Hodge
decomposition to singular varieties and provides a powerful tool for
studying the topology of algebraic varieties.

The theory of harmonic maps to non-positive curvature spaces (Eells--Sampson,
Uhlenbeck--Schoen--Yau) extends Kodaira's work to the geometric study
of harmonic maps between manifolds. The \textbf{Calabi--Yau theorem}---that
every K{\"a}hler class contains a unique Ricci-flat metric---uses the
same Hodge-theoretic framework. The Calabi--Yau theorem has important
applications in string theory, where Calabi--Yau manifolds serve as the
compactification spaces for extra dimensions.

Modern developments include the study of pluripotential theory on
complex manifolds, the theory of log canonical pairs, and the minimal
model program, all of which use Kodaira's vanishing and embedding
theorems as fundamental tools. The Kodaira--Spencer theory of
deformations underpins the modern study of moduli of curves, surfaces,
and higher-dimensional varieties. The deformation theory of complex
structures has found applications in mathematical physics, particularly
in the study of mirror symmetry and the moduli spaces of Calabi--Yau
manifolds.

The \textbf{Kodaira--Spencer map} $\rho: T_x \mathcal{M} \to H^1(X,
\Theta_X)$, where $\mathcal{M}$ is the moduli space of complex
structures on $X$ and $\Theta_X$ is the tangent sheaf, is the
fundamental tool for studying deformations. The Kodaira--Spencer
theorem states that $\rho$ is an isomorphism when $H^2(X, \Theta_X) =
0$, which implies that the moduli space is smooth and has dimension
$h^1(X, \Theta_X)$. When $H^2(X, \Theta_X) \neq 0$, the moduli space
may have singularities, and the Kodaira--Spencer map provides
information about the local structure near singular points.

The \textbf{Kodaira dimension} $\kappa(X)$, defined as the growth rate
of the plurigenera $P_m(X) = h^0(X, K_X^m)$, is a fundamental
birational invariant. The classification of surfaces by Enriques and
Kodaira, based on the Kodaira dimension, was the prototype for the
minimal model program in higher dimensions. Surfaces with $\kappa =
-\infty$ are rational or ruled, $\kappa = 0$ are K3, Enriques, Abelian,
or hyperelliptic, $\kappa = 1$ are properly elliptic, and $\kappa = 2$
are of general type. The Kodaira dimension is a birational invariant,
meaning that $\kappa(X) = \kappa(X')$ if $X$ and $X'$ are birationally
equivalent.

The \textbf{Barlet--Kodaira--Nakano} vanishing theorem for vector
bundles, proved by Nakano, extends the Kodaira vanishing theorem to
holomorphic vector bundles $E$ with positive curvature. The Nakano
vanishing theorem states: $H^q(X, K_X \otimes E) = 0$ for $q > 0$
when $E$ is positive, which has applications to the classification of
vector bundles on algebraic varieties. The Nakano positivity condition
is stronger than the Griffiths positivity condition and implies that
the curvature tensor has positive eigenvalues in all directions.

The \textbf{Bogomolov--Miyaoka--Yau inequality} $c_1^2 \leq 3c_2$ for
surfaces of general type, proved by Miyaoka (1973) and Yau (1977), is
a consequence of the Kodaira vanishing theorem and the theory of stable
vector bundles. The inequality provides a fundamental constraint on the
topology of surfaces of general type. Equality holds if and only if the
surface is a quotient of the unit ball in $\mathbb{C}^2$, which
connects the inequality to the theory of uniformization and hyperbolic
geometry.

The \textbf{Castelnuovo--de Franchis theorem} states that a compact
K{\"a}hler surface with two holomorphic 1-forms $\alpha, \beta$
satisfying $\alpha \wedge \beta = 0$ is fibered over a curve. This
theorem, which uses the Hodge decomposition and the Kodaira vanishing
theorem, is fundamental for the classification of surfaces. The theorem
provides a criterion for when a surface admits a fibration, and it has
been generalized to higher dimensions.

The \textbf{Bombieri--Enriques classification} of surfaces, completed
in the 1970s, uses the Kodaira dimension and the Hodge numbers to
classify all algebraic surfaces. The classification includes rational
surfaces ($\kappa = -\infty$), ruled surfaces ($\kappa = -\infty$),
surfaces with $\kappa = 0$ (K3, Enriques, bielliptic, Abelian),
surfaces with $\kappa = 1$ (properly elliptic), and surfaces of general
type ($\kappa = 2$). The classification is a landmark achievement of
twentieth-century algebraic geometry and demonstrates the power of
Kodaira's methods.

Kodaira's work on \textbf{compact K{\"a}hler manifolds} influenced the
development of \textbf{differential geometry} and \textbf{geometric
analysis}. The Calabi--Yau theorem, proved by Yau (1976), states that
every K{\"a}hler class on a compact K{\"a}hler manifold with $c_1 = 0$
contains a unique Ricci-flat metric. This result, which uses the same
Hodge-theoretic framework as Kodaira's work, is fundamental for string
theory and mirror symmetry. The Calabi--Yau theorem resolved the Calabi
conjecture, which had been open since 1954.

The \textbf{Kodaira--Spencer theory} has found applications far beyond
its original context. The deformation theory of complex structures
provides the foundation for the study of moduli spaces of algebraic
varieties, including the moduli space $\mathcal{M}_g$ of curves of
genus $g$. The local structure of $\mathcal{M}_g$ is governed by the
Kodaira--Spencer map, and the dimension of the Zariski tangent space at
a point $[C]$ is $h^1(C, \Theta_C) = 3g-3$. The global structure of
$\mathcal{M}_g$ has been studied extensively, and the boundary of the
moduli space (in various compactifications) encodes important geometric
information about stable curves.

The \textbf{Futaki invariant} and the theory of K{\"a}hler--Einstein
metrics extend Kodaira's ideas to the study of canonical metrics on
Fano manifolds. The existence of K{\"a}hler--Einstein metrics is
related to the stability of the underlying variety, a connection that
has led to deep results in the minimal model program and the
classification of Fano manifolds. The Yau--Tian--Donaldson conjecture
relates the existence of K{\"a}hler--Einstein metrics to K-stability,
a concept from geometric invariant theory.

Kodaira's techniques have been applied to the study of
\textbf{hyperk{\"a}hler manifolds}, which are compact K{\"a}hler
manifolds with a rich structure of complex structures. The moduli
spaces of hyperk{\"a}hler manifolds, including K3 surfaces and their
higher-dimensional analogues, have been extensively studied using
Kodaira's methods. The Bogomolov--Miyaoka--Yau inequality and the
classification of surfaces of general type both rely on the vanishing
theorems developed by Kodaira. Hyperk{\"a}hler manifolds are important
in mathematical physics, particularly in the study of supersymmetric
gauge theories and string compactifications.

The theory of \textbf{mixed Hodge structures}, developed by Deligne,
extends Kodaira's Hodge decomposition to singular varieties and
provides a powerful tool for studying the topology of algebraic
varieties. The mixed Hodge structure on the cohomology of a smooth
variety $X$ reflects the geometric structure of $X$ and its
compactifications. The weight filtration on the cohomology measures the
complexity of the singularities, and the monodromy action encodes
information about the topology of the variety near singular fibers.

The \textbf{mirror symmetry} conjecture, which relates the geometry of
Calabi--Yau manifolds, relies heavily on Kodaira's work. The Hodge
numbers of a Calabi--Yau threefold and its mirror are related by the
symmetry $h^{1,1}(X) = h^{2,1}(X')$ and $h^{2,1}(X) = h^{1,1}(X')$.
This symmetry has been verified for many examples and has led to deep
results in enumerative geometry, including the computation of
Gromov--Witten invariants. The study of variations of Hodge structure
on the moduli space of Calabi--Yau manifolds is a central topic in
mirror symmetry.

\begin{center}\rule{0.5\linewidth}{0.5pt}\end{center}

\section*{Summary}
\addcontentsline{toc}{section}{Summary}

Kodaira developed the theory of harmonic integrals on K{\"a}hler
manifolds, proving the Hodge decomposition and the Kodaira vanishing
theorem. His embedding theorem characterized projective manifolds, and
his deformation theory founded the modern study of moduli spaces of
complex varieties. Serre duality, extended to compact complex manifolds
by Kodaira, provides a fundamental relationship between cohomology
groups and their duals. The Kodaira dimension and the classification of
algebraic surfaces, based on the growth of plurigenera, remain central
to the minimal model program. Kodaira's work unified complex geometry,
differential geometry, and algebraic geometry, creating a framework
that continues to drive research in mathematics.
